\begin{center}
	\zihao{4}\heiti
	% 中国科学院大学\\
	\LARGE{中国科学院研究生院\\$\boldsymbol{2011}$年招收硕士研究生入学统一考试试卷\\
	科目名称：电动力学（回忆版）}

\end{center}
\noindent
考生须知：\\
1.本试卷满分为150分,全部考试时间总计180分钟。\\
2.所有答案必须写在答题纸上,写在试题纸上或草稿纸上一律无效。

\vspace*{2ex}
\hrule
\vspace*{4ex}

\fancypagestyle{mypagestyle}{%
	\fancyhf{}% Clear header/footer
	\fancyfoot[OR]{第~\thepage~页，共2页}%
	\fancyfoot[EL]{第~\thepage~页，共2页}%
	\fancyfoot[OL]{科目名称：电动力学 }%
	\fancyfoot[ER]{科目名称：电动力学 }%
	\renewcommand{\headrulewidth}{0pt}% Header rule of .4pt
	\renewcommand{\footrulewidth}{0.5pt}% Header rule of .4pt
}
\fancypagestyle{plain}{%
	\fancyhf{}% Clear header/footer
	\fancyfoot[OR]{第~\thepage~页，共2页}%
	\fancyfoot[EL]{第~\thepage~页，共2页}%
	\fancyfoot[OL]{科目名称：电动力学 }%
	\fancyfoot[ER]{科目名称：电动力学 }%
}

\setcounter{page}{1}%设置页面计数器
\pagestyle{mypagestyle}

\onehalfspacing
% 排版结束，下面是题目
\vspace{0.4cm}

\noindent
一、简答题
\vspace{0.2cm}

\noindent
1.\hspace{0.8em}写出介质中有源麦克斯韦方程组的微分形式及积分形式。
\vspace{0.2cm}

\noindent
2.\hspace{0.8em}写出自由电荷受到的力密度表达式，试写出电磁场的边界关系
\vspace{0.2cm}

\noindent
3.\hspace{0.8em}试推导各项同性均匀介质内的时谐波的亥姆霍兹方程
\vspace{0.2cm}

\noindent
4.\hspace{0.8em}试写出谐振腔的本征频率表达式
\vspace{0.2cm}

\noindent
二、\hspace{0.2em}单项选择题
\vspace{0.2cm}

\noindent
1.\hspace{0.8em}均匀带电的椭球面带总电量Q，坐标原点在椭球中心，其
\vspace{0.2cm}

A.\hspace{0.8em}电偶极矩和电四极矩都不为零 \hspace{2.6em} B.\hspace{0.8em}电偶极矩为零，电四极矩不为零
\vspace{0.2cm}

C.\hspace{0.8em}电偶极矩不为零，电四极矩为零\hspace{2.1em} D.\hspace{0.8em}电偶极矩，电四极矩都为零
\vspace{0.2cm}

\noindent
2.\hspace{0.8em}中科大09年电动A选择第三题）电偶极辐射在远处的能流密度随距离$r$的变化关系为\vspace{0.2cm}

A.\hspace{0.8em}正比于$r$ \hspace{2em} B.\hspace{0.8em}与$r$无关 \hspace{2em} C.\hspace{0.8em}正比于$r^{-1}$\hspace{2em} D.\hspace{0.8em}正比于$r^{-2}$
\vspace{0.2cm}

\noindent
3.\hspace{0.8em}中科大09年电动A选择第四题）静止$\mu$子的平均寿命是$2.2 \times 10^{-6} s $。在实验室中，从高能加速器出来的$\mu$子以0.6c运动。在实验室中观察，$\mu$子衰变前的平均飞行距离为;
\vspace{0.2cm}

A.\hspace{0.8em} 825 m \hspace{2em} B.\hspace{0.8em} 316.8 m \hspace{2em} C.\hspace{0.8em} 396 m\hspace{2em} D.\hspace{0.8em} 495 m
\vspace{0.2cm}

\noindent
4.\hspace{0.8em}中科大10年电动A选择第五题）设区域V内给定自由电荷分布$\rho(\Vec{x})$，在V的边界S上给定什么条件能保证V内的静电场唯一地确定
\vspace{0.2cm}

A.\hspace{0.8em}电荷分布$\sigma|_s$ \hspace{5em}
B.\hspace{0.8em}电势$\varphi|_s$以及电势的法向导数$\displaystyle{ \frac{\partial \varphi}{\partial n}}|_s$\vspace{0.1cm}

C.\hspace{0.8em}电势的切线导数\hspace{4.2em}
D.\hspace{0.8em}电势$\varphi|_s$或电势的法向导数$\displaystyle{ \frac{\partial \varphi}{\partial n}}|_s$ \vspace{0.2cm}

\noindent
5.\hspace{0.8em}中科大08年电动A选择第二题）频率为100 MHz的电磁波对铜的穿透深度为$\delta \sim 0.7 \times 10^{-3}$cm。当电磁波的频率为1 MHz时，铜的穿透深度为？
\vspace{0.2cm}

A.\hspace{0.8em} $0.7 \times 10^{-1} $cm \hspace{2em}
B.\hspace{0.8em} $0.7 \times 10^{-2} $cm \hspace{2em}
C.\hspace{0.8em} $0.7 \times 10^{-3} $cm \hspace{2em}
D.\hspace{0.8em} $0.7 \times 10^{-4} $cm 

\vspace{0.4cm}

\noindent
三、\hspace{0.2em}有一个均匀带电的薄导体壳，半径为$R_0$，总电荷量为$q$，今使球壳绕自身某一直径以角速度$\omega$转动，求球内外的磁场$B$。（课后题3.13）\vspace{1cm}

\noindent
四、\hspace{0.2em}某一均匀磁场中带电粒子e做半径为a的非相对论性圆周运动，回旋频率为$\omega$，求远处的辐射电磁场，辐射能流和辐射功率。(与课后题5.11相似)
